\section*{使用说明}
本讲义面向已经学过变量、分支、循环、数组和函数，希望系统训练 CCF-CSP 认证前三题的学生。学习主线是：先把短程序写准确，再学习减少重复计算的方法，最后把长题意拆成可靠的程序模块。

正文中的教学例题均给出完整的任务、输入输出、规模限制、样例、解答和参考实现，可以脱离外部网站学习。例题是围绕原资料知识点重新组织的教学题，\textbf{不等同于所链接平台题目的完整原版题面}。原题链接用于进一步练习，不能直接把教学例题的代码提交到规则不同的原题。

每章建议按三个步骤学习：先自己推演例题，再解释参考代码中的关键不变量，最后遮住答案完成章末练习。不会做时优先寻找“状态究竟表示什么”“哪一步重复计算”“哪个边界没有定义”，而不是立即记忆模板。

参考程序使用 C++17。为减少重复，使用竞赛环境常见的 \texttt{bits/stdc++.h}；该头文件属于 GNU 实现。其他编译器可换成程序实际使用的标准头文件，如 iostream、vector、algorithm、map、set、queue、sstream、string、limits、stdexcept。学习时使用 GNU C++17 编译即可，具体认证提交环境应在考试前核对。

第1--10章构成基础算法主线，第11章包含五个完整模拟案例，第12章提供带答案的综合训练。附录补充多重背包、多维背包、背包撤销等原资料中的提高内容。高级内容不影响先完成基础与模拟主线。

\important{一个始终要回答的问题}
一段代码通过样例，只说明它在那个输入上表现正确。每次写完都应检查：最小规模、最大规模、重复值、边界位置、不可达状态，以及题面允许但样例没有出现的情况。
\tableofcontents
\newpage
